\documentclass[hyperref={breaklinks=true},fleqn,mathserif]{beamer}

\usepackage{graphicx}
\usepackage{booktabs}
\usepackage{array}
\usepackage{amssymb}
\usepackage{listings}
\usepackage{tikz}
\usetikzlibrary{arrows.meta,positioning,fit,calc,shapes.geometric}

%%%%%%%%%%%%%%%%%%%%%%%%%%%%%%%%%%%%%%%%%%%%%%%%%%%%%%%%%%%%%%%%%%%%%%%%%%%%%%%%
\usepackage[utf8x]{inputenc}
\usepackage[L7x]{fontenc}
\usepackage[english,lithuanian]{babel}
%%%%%%%%%%%%%%%%%%%%%%%%%%%%%%%%%%%%%%%%%%%%%%%%%%%%%%%%%%%%%%%%%%%%%%%%%%%%%%%%

\usetheme{Warsaw}
\setbeamertemplate{navigation symbols}{}
\setbeamersize{text margin left=1em,text margin right=1em}

\title[Dinaminė žmogžudysčių duomenų bazė]{Dinaminė žmogžudysčių duomenų bazė}
\subtitle{Rezultatai, algoritmai ir tolesni planai}
\author[\textit{et al.}]{\textbf{Paulius Jankauskas}, Tomas Sarafin ir \\ Justas Daugirdas}
\institute{Vilniaus universitetas \\ Matematikos ir informatikos fakultetas }
\date{2026 m. spalio 21 d.}

\defbeamertemplate*{footline}{shadow theme}
{%
    \leavevmode%
    \hbox{%
        \begin{beamercolorbox}[wd=.5\paperwidth,ht=2.5ex,dp=1.125ex,leftskip=.3cm plus1fil,rightskip=.3cm]{author in head/foot}%
            \usebeamerfont{author in head/foot}\insertframenumber\,/\,\inserttotalframenumber\hfill\insertshortauthor
        \end{beamercolorbox}%
        \begin{beamercolorbox}[wd=.5\paperwidth,ht=2.5ex,dp=1.125ex,leftskip=.3cm,rightskip=.3cm plus1fil]{title in head/foot}%
            \usebeamerfont{title in head/foot}\insertshorttitle%
        \end{beamercolorbox}%
    }%
    \vskip0pt%
}

% Code listings (code itself is ASCII only). Computer Modern has no L7x
% typewriter font, so use Latin Modern's for code and \texttt.
\renewcommand{\ttdefault}{lmtt}
\definecolor{codebg}{RGB}{245,245,250}
\definecolor{codekw}{RGB}{38,38,134}
\definecolor{codecm}{RGB}{110,110,110}
\definecolor{codestr}{RGB}{150,40,40}
\definecolor{termbg}{RGB}{30,32,36}
\definecolor{termfg}{RGB}{225,225,225}
\definecolor{termhi}{RGB}{120,210,120}
\definecolor{termerr}{RGB}{240,120,110}
\lstset{
    language=Python,
    basicstyle=\ttfamily\scriptsize,
    keywordstyle=\color{codekw}\bfseries,
    commentstyle=\color{codecm}\itshape,
    stringstyle=\color{codestr},
    backgroundcolor=\color{codebg},
    frame=single, rulecolor=\color{codebg},
    columns=fullflexible, keepspaces=true,
    showstringspaces=false, breaklines=true,
    xleftmargin=2pt, xrightmargin=2pt,
}
\lstdefinestyle{sql}{language=SQL, morekeywords={ROLE,LOGIN,NOINHERIT,LIMIT,REVOKE}}
\lstdefinestyle{term}{
    language={}, basicstyle=\ttfamily\tiny\color{termfg},
    backgroundcolor=\color{termbg}, rulecolor=\color{termbg},
    columns=fullflexible, breaklines=false,
    emph={detective}, emphstyle=\color{termhi}\bfseries,
    emph={[2]rejected,ERROR,wrong}, emphstyle={[2]\color{termerr}\bfseries},
}

% TikZ styles
\tikzset{
    box/.style={draw=structure, fill=structure!8, rounded corners=3pt, align=center,
                font=\scriptsize, inner sep=4pt, minimum height=8mm},
    hot/.style={box, draw=red!70!black, fill=red!8},
    good/.style={box, draw=green!45!black, fill=green!8},
    arr/.style={-{Latex[length=2mm]}, thick, structure},
    lbl/.style={font=\tiny, fill=white, inner sep=1pt},
    tbl/.style={draw=structure, fill=white, rounded corners=2pt, align=left,
                font=\tiny, execute at begin node={\hyphenpenalty=10000\relax},
                inner sep=2.5pt, text width=32.5mm},
    fk/.style={-{Latex[length=1.6mm]}, structure!70},
}
\newcommand{\tname}[1]{\textbf{\scriptsize #1}\\[1pt]}
\newcommand{\yes}{\textcolor{green!50!black}{$\checkmark$}}
\newcommand{\no}{\textcolor{red!70!black}{$\times$}}

\begin{document}

\frame[plain]{\titlepage}

% ============================================================
% 1. KĄ PADARĖME
% ============================================================
\begin{frame}{Užduotis ir ką padarėme}

\begin{columns}[T]
\begin{column}{0.46\textwidth}
\begin{block}{Užduotis}
\footnotesize
Policijos duomenų bazėje yra daug nusikaltimų bylų. Reikia sistemos, kuri iš šių duomenų
\textbf{kiekvienam žaidėjui sukuria naują žmogžudystės bylą}: kitą auką, žudiką, liudininkus,
faktus ir ryšius. Žaidėjas bylą tiria \textbf{SQL užklausomis}.
\par\medskip
Pavyzdys: \texttt{ctf.cs.vu.lt/detective}, bet ten atsakymas visiems vienodas.
\end{block}
\end{column}

\begin{column}{0.52\textwidth}
\begin{exampleblock}{Kas jau veikia}
\footnotesize
\begin{itemize}
    \setlength{\itemsep}{2pt}
    \item Bazinė \textbf{policijos duomenų bazė} (PostgreSQL).
    \item \textbf{Scenarijų variklis}: iš jos sukuria naują bylą.
    \item Ta pati \textbf{sėkla (seed)} atkuria tą pačią bylą.
    \item Kiekviena byla turi \textbf{vienintelį} sprendinį (patikrinta SQL).
    \item 3 sunkumo lygiai, užuominos, taškai.
    \item Žaidimas \textbf{terminale}, su apribota DB role.
\end{itemize}
\end{exampleblock}

\begin{block}{Technologijos}
\footnotesize Python 3, PostgreSQL 16, psycopg 3, virtualios mašinos (Linux).
\textit{Planuojama:} Faker duomenims generuoti.
\end{block}
\end{column}
\end{columns}

\end{frame}

% ============================================================
% 2. ARCHITEKTŪRA (VM)
% ============================================================
\begin{frame}{Sistemos architektūra: virtualios mašinos}

\begin{center}
\begin{tikzpicture}
    % Player side
    \node[box, minimum width=22mm] (pl) at (0,0.9) {Žaidėjas\\\tiny terminalas};
    \node[box, minimum width=22mm] (ps) at (0,-1.05) {\texttt{psql}\\\tiny (nebūtina)};

    % Program
    \node[box, minimum width=30mm] (eng) at (4.0,2.25) {Scenarijų variklis\\\tiny \texttt{scenario.py}};
    \node[box, minimum width=30mm] (cli) at (4.0,0.9) {\texttt{detective.py}\\\tiny new, play, hint, accuse};
    \node[draw=structure, dashed, rounded corners, fit=(eng)(cli), inner sep=4pt,
          label={[font=\tiny, text=structure, inner sep=1pt]above:Python}] (py) {};

    % PostgreSQL
    \node[good, minimum width=31mm] (pol) at (8.4,2.25) {\texttt{police}\\\tiny bazinis archyvas};
    \node[good, minimum width=31mm] (game) at (8.4,0.9) {\texttt{game}\\\tiny atsakymų maišos, užuominos};
    \node[hot, minimum width=31mm] (case) at (8.4,-0.9) {\texttt{case\_<id>}\\\tiny po vieną kiekvienai bylai};
    \node[draw=structure, dashed, rounded corners, fit=(pol)(game)(case), inner sep=4pt,
          label={[font=\tiny, text=structure, inner sep=1pt]above:PostgreSQL 16}] (pg) {};

    \node[draw=structure, very thick, rounded corners, fit=(py)(pg), inner sep=4pt, inner ysep=10pt,
          label={[font=\scriptsize\bfseries, text=structure]below:Serverio VM (Linux)}] (vm) {};

    \draw[arr] (pl) -- node[lbl, above, pos=0.35] {SSH} (cli);
    \draw[arr] (eng) -- node[lbl, right] {byla} (cli);
    \draw[arr] (eng) -- node[lbl, above] {skaito} (pol);
    \draw[arr] (cli) -- node[lbl, above] {admin} (game);
    \draw[arr] ($(cli.east)+(0,-0.2)$) -- ++(0.55,0) |- node[lbl, pos=0.3, right] {sukuria} ($(case.west)+(0,0.15)$);
    \draw[arr, red!70!black] (ps.east |- 0,-1.05) -- node[lbl, above, pos=0.35] {rolė \texttt{p\_<id>}: tik SELECT} (case.west |- 0,-1.05);
\end{tikzpicture}
\end{center}

\vspace{-1mm}
{\footnotesize
\begin{itemize}
    \item Be API ir Docker: programa ir duomenų bazė veikia virtualioje mašinoje, žaidėjas jungiasi per SSH.
    \item Žaidėjo SQL \textbf{niekada} nevykdomas administratoriaus ryšiu, tik bylos role \texttt{p\_<id>}.
\end{itemize}
}

\end{frame}

% ============================================================
% 3. DUOMENŲ BAZĖ
% ============================================================
\begin{frame}{Duomenų bazė: policijos archyvas}

\begin{center}
\begin{tikzpicture}
    \node[tbl] (dl) at (0,1.9) {\tname{drivers\_license}\underline{id}, gender, height, hair\_color, \textbf{plate\_number}, car\_make, car\_model};
    \node[tbl, draw=red!70!black, very thick] (p) at (4.05,1.9) {\tname{person}\underline{id}, \textbf{name}, \textit{license\_id}, address\_number, \textbf{address\_street\_name}};
    \node[tbl] (csr) at (8.1,1.9) {\tname{crime\_scene\_report}date, type, city, description\\[1pt]\textit{be ryšių: įkalčiai tekste}};
    \node[tbl] (int) at (0,-0.15) {\tname{interview}\textit{person\_id}, transcript};
    \node[tbl] (fb) at (4.05,-0.15) {\tname{facebook\_event\_checkin}\textit{person\_id}, event\_name, date};
    \node[tbl] (gm) at (8.1,-0.15) {\tname{get\_fit\_now\_member}\underline{id}, \textit{person\_id}, membership\_status};
    \node[tbl] (gc) at (8.1,-1.75) {\tname{get\_fit\_now\_check\_in}\textit{membership\_id}, check\_in\_date};

    \draw[fk] (p) -- (dl);
    \draw[fk] (int) -- (p);
    \draw[fk] (fb) -- (p);
    \draw[fk] (gm) -- (p);
    \draw[fk] (gc) -- (gm);
\end{tikzpicture}
\end{center}

\vspace{-1mm}
{\footnotesize
\begin{itemize}
    \item \texttt{police}: bazinis rinkinys (\texttt{police\_data.sql}): 500 asmenų, 200 ataskaitų. Jie bylose \textbf{perpanaudojami}.
    \item Sumažinta: 7 lentelės, tik žaidimui reikalingi stulpeliai (be \texttt{income}, \texttt{ssn} ir pan.).
    \item Kiekviena byla gauna \textbf{kopiją} su ta pačia struktūra savo schemoje \texttt{case\_<id>}.
\end{itemize}
}

\end{frame}

% ============================================================
% 4. ĮKALČIŲ GRANDINĖ
% ============================================================
\begin{frame}{Bylos generavimas: įkalčių grandinė}

\begin{center}
\begin{tikzpicture}[node distance=4mm and 4mm]
    \node[box, text width=17mm] (r) {Nusikaltimo\\ataskaita};
    \node[box, right=of r, text width=17mm] (w) {Liudininkai\\\tiny gatvė, vardas};
    \node[box, right=of w, text width=17mm] (i) {Apklausos};
    \node[box, right=of i, text width=22mm] (g) {Klubo narystė\\+ numerio dalis};

    \node[box, below=7mm of g, text width=22mm] (c) {Registracija\\klube tą dieną};
    \node[hot, left=of c, text width=17mm] (k) {\textbf{Žudikas}};
    \node[box, left=of k, text width=17mm] (ki) {Žudiko\\apklausa};
    \node[hot, left=of ki, text width=17mm] (m) {\textbf{Užsakovas}\\\tiny renginiai, auto};

    \draw[arr] (r) -- (w); \draw[arr] (w) -- (i); \draw[arr] (i) -- (g);
    \draw[arr] (g) -- (c); \draw[arr] (c) -- (k); \draw[arr] (k) -- (ki); \draw[arr] (ki) -- (m);
    \node[font=\tiny, text=structure, below=0.5mm of ki, xshift=-10mm] {tik Vidutinė / Sunki};
\end{tikzpicture}
\end{center}

\vspace{-2mm}
\begin{columns}[T]
\begin{column}{0.43\textwidth}
\begin{block}{Principas}
\footnotesize
\begin{enumerate}
    \item Auka, liudininkai ir įtariamieji \textbf{parenkami iš archyvo}, o jų faktai pakeičiami.
    \item Pirmiausia parenkamas \textbf{nusikaltėlis}, įrodymai išvedami iš jo.
    \item Kiekvienas klaidinantis įtariamasis neatitinka \textbf{lygiai vieno} kriterijaus.
\end{enumerate}
\end{block}
\end{column}
\begin{column}{0.55\textwidth}
\scriptsize
\setlength{\tabcolsep}{3pt}
\begin{tabular}{@{}lcccc@{}}
\toprule
 & Klubo & Auksinė & Numeris & Buvo \\
 & narys & narystė & sutampa & klube \\
\midrule
\textbf{Žudikas} & \yes & \yes & \yes & \yes \\
Apgaulė 1 & \yes & \yes & \no & \yes \\
Apgaulė 2 & \yes & \no & \yes & \yes \\
Apgaulė 3 & \yes & \yes & \yes & \no \\
Apgaulė 4 & \no & -- & \yes & -- \\
\bottomrule
\end{tabular}
\par\smallskip
{\scriptsize Sunkesniame lygyje: daugiau apgaulių, trumpesnė numerio dalis.}
\end{column}
\end{columns}

\end{frame}

% ============================================================
% 5. ALGORITMAS
% ============================================================
\begin{frame}{Generavimo algoritmas}

\begin{center}
\begin{tikzpicture}[node distance=5mm]
    \node[box, text width=19mm] (s) at (0,1.6) {Sėkla + lygis};
    \node[box, text width=19mm] (b) at (2.9,1.6) {Bazinio archyvo kopija};
    \node[box, text width=19mm] (x) at (5.8,1.6) {Pašalinti atsitiktinius sutapimus};
    \node[hot, text width=19mm] (k) at (8.7,1.6) {Parinkti žudiką \textbf{(pirmą)}};

    \node[box, text width=19mm] (e) at (8.7,-0.9) {Įrodymai, apgaulės, liudininkai, auka};
    \node[box, text width=19mm] (l) at (5.8,-0.9) {Įkelti į \texttt{case\_<id>}\\\tiny viena transakcija};
    \node[box, draw=orange!80!black, fill=orange!10, very thick, text width=19mm] (v) at (2.9,-0.9) {\textbf{SQL patikra}\\\tiny ar kiekvienas etapas turi lygiai 1 atsakymą?};
    \node[good, text width=19mm] (ok) at (0,-0.9) {Rolė \texttt{p\_<id>}, atsakymų maišos, COMMIT};

    \draw[arr] (s) -- (b); \draw[arr] (b) -- (x); \draw[arr] (x) -- (k);
    \draw[arr] (k) -- (e); \draw[arr] (e) -- (l); \draw[arr] (l) -- (v);
    \draw[arr, green!45!black] (v) -- node[lbl, above=1pt] {Taip} (ok);
    \draw[arr, red!70!black] (v) -- node[lbl, right, align=left] {Ne: ROLLBACK,\\bandymas + 1} (b);
\end{tikzpicture}
\end{center}

\vspace{-2mm}
{\footnotesize
\begin{itemize}
    \setlength{\itemsep}{2pt}
    \item Visi atsitiktiniai pasirinkimai: \texttt{random.Random(f"\{seed\}:\{attempt\}")}, todėl ta pati sėkla \(\rightarrow\) ta pati byla.
    \item Prieš kuriant bylą archyve pakeičiami atsitiktiniai atitikmenys: numeriai su ta pačia dalimi, to paties renginio lankymai.
    \item Patikra vykdo \textbf{tas pačias SQL užklausas}, kurias turės parašyti žaidėjas.
\end{itemize}
}

\end{frame}

% ============================================================
% 6. KODAS
% ============================================================
\begin{frame}[fragile]{Kodas: kontroliuojamas atsitiktinumas}

\vspace{-1mm}
\begin{lstlisting}[basicstyle=\ttfamily\fontsize{6.8}{8}\selectfont]
# scenario.py: the killer is cast first, from people already in the archive
killer, *decoys = rng.sample(drivers("male"), 1 + p["decoys"])
make_suspect(killer, "gold", gym, True)               # matches every clue
misses = [("gold", gym, False),                       # wrong plate
          ("silver", gym, True),                      # not gold
          ("gold", gym - dt.timedelta(days=1), True), # at the gym on another day
          (None, None, True)]                         # not a member
for i, d in enumerate(decoys):                        # each decoy misses exactly one
    make_suspect(d, *misses[i % len(misses)])

# detective.py: keep the case only if the player's SQL path has one answer
for attempt in range(20):
    case = build_case(base, seed, attempt, level)
    try:
        with conn.transaction():
            provision(conn, case, case_id, seed, level)  # case_<id>, VERIFY_SQL, role
        return case, attempt
    except NotUnique:                                    # rolled back: next attempt
        continue
\end{lstlisting}

\vspace{-2mm}
\begin{lstlisting}[style=term]
$ python detective.py selftest --count 100
easy    100 seeds: one answer per stage, same case every time, 0 retries, 199 ms/case
medium  100 seeds: one answer per stage, same case every time, 0 retries, 198 ms/case
hard    100 seeds: one answer per stage, same case every time, 0 retries, 192 ms/case
\end{lstlisting}

\end{frame}

% ============================================================
% 7. PRISIJUNGIMAS PRIE DB
% ============================================================
\begin{frame}[fragile]{Kaip jungiamės prie duomenų bazės}

\begin{center}
\begin{tikzpicture}
    \node[font=\scriptsize\bfseries, text=structure, anchor=west] at (-1.2,1.55) {Administratoriaus ryšys (\texttt{murder\_admin})};
    \node[font=\scriptsize\bfseries, text=red!70!black, anchor=west] at (-1.2,-0.35) {Žaidėjo ryšys (\texttt{p\_<id>})};

    \node[box, text width=21mm] (a1) at (0,0.85) {\textbf{1. new}\\\tiny skaito \texttt{police}};
    \node[box, text width=21mm] (a2) at (2.75,0.85) {\textbf{2.} sukuria \texttt{case\_<id>} ir rolę \texttt{p\_<id>}};
    \node[box, text width=21mm] (a3) at (5.5,0.85) {\textbf{3.} slaptažodis ir atsakymų maišos \(\rightarrow\) \texttt{game}};
    \node[box, text width=21mm] (a4) at (8.25,0.85) {\textbf{4. play}\\\tiny paima slaptažodį};

    \node[hot, text width=21mm] (p1) at (8.25,-1.05) {\textbf{5.} prisijungia kaip \texttt{p\_<id>}};
    \node[hot, text width=21mm] (p2) at (5.5,-1.05) {\textbf{6.} \texttt{SELECT} tik iš \texttt{case\_<id>}};
    \node[good, text width=21mm] (p3) at (2.75,-1.05) {\textbf{7.} \texttt{hint} / \texttt{accuse}\\\tiny vėl per admin};

    \draw[arr] (a1) -- (a2); \draw[arr] (a2) -- (a3); \draw[arr] (a3) -- (a4);
    \draw[arr, red!70!black] (a4) -- (p1); \draw[arr, red!70!black] (p1) -- (p2);
    \draw[arr] (p2) -- (p3);
    \draw[dashed, structure!50] (-1.2,-0.15) -- (9.6,-0.15);
\end{tikzpicture}
\end{center}

\vspace{-2mm}
\begin{columns}[T]
\begin{column}{0.58\textwidth}
\begin{lstlisting}[basicstyle=\ttfamily\fontsize{6.2}{7.4}\selectfont]
def player_connection(case_id, password):
    # same server and database, but another user
    info = conninfo_to_dict(ADMIN_DSN) | {
        "user": f"p_{case_id}", "password": password,
        "options": "-c default_transaction_read_only=on"
                   " -c statement_timeout=3000"}
    return psycopg.connect(make_conninfo(**info))

wrapped = f"SELECT * FROM (\n{query}\n) LIMIT 200"
player.execute(wrapped, prepare=True)
\end{lstlisting}
\end{column}
\begin{column}{0.40\textwidth}
\begin{lstlisting}[style=term]
detective> DELETE FROM person;
rejected: only SELECT queries are allowed
detective> SELECT * FROM game.cases;
ERROR: permission denied for schema game
detective> SELECT * FROM police.person;
ERROR: permission denied for schema police
detective> SELECT pg_sleep(10);
ERROR: canceling statement due to
       statement timeout
\end{lstlisting}
\end{column}
\end{columns}

\end{frame}

% ============================================================
% 8. DEMONSTRACIJA TERMINALE
% ============================================================
\begin{frame}[fragile]{Kaip tai veikia terminale}

\begin{columns}[T]
\begin{column}{0.68\textwidth}
\begin{lstlisting}[style=term]
$ python detective.py new medium --seed 42
case 7e979be7  level=medium  seed=42  attempt=0  (249 ms)
schema case_7e979be7, role p_7e979be7 (read-only, 3000 ms limit)

$ python detective.py play 7e979be7
detective> SELECT description FROM crime_scene_report
        -> WHERE type='murder' AND city='SQL City' AND date=20260701;
 Heather Lee was found dead. There were 2 witnesses. The
 first lives at the last house on "Union Street". The second,
 named Deborah, lives somewhere on "Willow Way".
detective> SELECT p.name, m.membership_status AS status,
        ->   c.check_in_date AS day, d.plate_number AS plate
        -> FROM get_fit_now_member m JOIN ... WHERE m.id LIKE '95U%';
 name           | status | day      | plate
----------------+--------+----------+--------
 Brandon Torres | silver | 20260626 | 167QJ3
 Brandon Kim    | gold   | 20260625 | 97QJ3X
 Andrew Ramirez | gold   | 20260626 | K7QJ3M
 Donald Taylor  | gold   | 20260626 | PYXT4S
detective> accuse Donald Taylor
wrong killer. Keep digging.
detective> accuse Andrew Ramirez
Correct, that's the killer! But someone hired them. ...
\end{lstlisting}
\end{column}

\begin{column}{0.30\textwidth}
\footnotesize
\begin{enumerate}
    \setlength{\itemsep}{2pt}
    \item Sukuriama byla: nauja schema ir rolė.
    \item Ataskaita: auka ir 2 liudininkai.
    \item Liudininkai: narystė \texttt{95U...}, auksinė, numeryje \texttt{7QJ3}, birželio 26 d.
    \item \textbf{Torres}: ne auksinė. \textbf{Kim}: kita diena. \textbf{Taylor}: kitas numeris.
    \item Žudikas rastas, toliau ieškoma užsakovo.
\end{enumerate}
\end{column}
\end{columns}

\end{frame}

% ============================================================
% 9. IŠŠŪKIAI IR ATEITIES PLANAI
% ============================================================
\begin{frame}{Iššūkiai ir ateities planai}

\begin{columns}[T]
\begin{column}{0.58\textwidth}
\scriptsize
\begin{tabular}{@{}>{\raggedright\arraybackslash}p{0.44\textwidth}>{\raggedright\arraybackslash}p{0.52\textwidth}@{}}
\toprule
\textbf{Problema} & \textbf{Sprendimas} \\
\midrule
Byla gali turėti kelis atsakymus & Nusikaltėlis pirmas, SQL patikra, kitas bandymas \\ \addlinespace[3pt]
Archyvo duomenys atsitiktinai atitinka įkalčius & Sutapimai pašalinami prieš kuriant bylą \\ \addlinespace[3pt]
Bendra rolė + \texttt{SET ROLE} leistų patekti į kitą bylą & Atskira prisijungimo rolė kiekvienai bylai \\ \addlinespace[3pt]
Kelios komandos vienoje užklausoje & SQL filtras ir paruoštas sakinys \\ \addlinespace[3pt]
\texttt{CROSS JOIN}, \texttt{pg\_sleep} apkrautų serverį & 3 s laiko limitas ir \texttt{LIMIT} \\ \addlinespace[3pt]
Atsakymai neturi būti matomi & \texttt{sha256(druska:vardas)} schemoje \texttt{game} \\
\bottomrule
\end{tabular}
\end{column}

\begin{column}{0.40\textwidth}
\begin{block}{Toliau}
\footnotesize
\begin{itemize}
    \item \textbf{Faker}: realistiškesni ir didesni duomenų rinkiniai.
    \item Daugiau bylų šablonų ir motyvų.
    \item Atskiros VM duomenų bazei ir žaidėjams.
    \item Automatinis senų bylų šalinimas.
    \item Rezultatų lentelė.
\end{itemize}
\end{block}
\end{column}
\end{columns}

\vspace{3mm}
\begin{center}
\large Ačiū už dėmesį! Klausimai?
\end{center}

\end{frame}

\end{document}
